\interlude{至福一瞬}
\footnotetext{%
    译注：标题化自歌德《浮士德》中的名句：“如果我对某个瞬间说：停留一下吧，你多么美呀！”（绿原先生译）。此处指一个幸福圆满到令人愿其永远停驻的瞬间。
}
若要深入讨论《魔法少女小圆》，我们几乎不可能绕开歌德的《浮士德》。在动画第二集中，墙上的德语涂鸦便直接引自这部作品。这个关于人与魔鬼订立契约的古老传说，经歌德重新讲述之后，构成了贯穿全剧的重要思想底本。

先简单概括一下《浮士德》的故事：
浮士德\footnote{%
    译注：
    《浮士德》中出场人物的译名，均参考人民文学出版社“名著名译丛书”《浮士德》绿原先生的译本。
}是一位博学睿智的老人，却无法从自己的人生中获得任何欢愉。
为了重获青春，换一种方式重新生活，他与恶魔梅菲斯特订立了契约。
梅菲斯特答应让他尽享过去错失的一切快乐。
但作为代价，一旦浮士德抵达至福之境，幸福得甚至希望时间就此停驻、让那一刻永远延续，他便会立即死去，堕入地狱。

第一部初版于1808年，后于1828年修订，主要讲述浮士德追求年轻女子玛格丽特的故事。
玛格丽特也常被称作格蕾琴，那是她的昵称。
浮士德杀死她的兄长之后，一度离去，前往参加“瓦尔普吉斯之夜”的狂欢。
按照德国民间传说，每逢此夜，女巫与恶魔都会聚集在布罗肯山上纵欲作乐。
浮士德归来时，格蕾琴已经神志失常，身陷囹圄。
她曾生下浮士德的孩子，但孩子被夺走了\footnote{%
   译注：原文写道 ``..., and she gave birth to his child but it was taken away.''
   直译为：她生下了他的孩子，但孩子却被带走／夺去了。
   这里有些歧义。
   实际上，在《浮士德》中，我们从身陷地牢的玛格丽特口中得知，是她亲手溺死了婴儿。
}。
浮士德试图救她出狱，然而格蕾琴却深陷幻乱，根本无法理解眼前发生的一切。
卫兵将至，浮士德只得将她抛下，独自逃亡\footnote{%
    本节若无另行说明，凡涉及《浮士德》情节之处，均以贝阿德·泰勒的英译本为据\cite{goethe2005faust}。
}。


第二部出版于1832年，也就是歌德去世的那一年，情节则奇异得多。
此时的浮士德再次老去，却已功成名就、富甲一方，并成为了一位法力强大的术士。
他穿梭时空，与特洛伊的海伦有过一段情缘，又借助一支恶魔军队赢得战争。
到了故事终章，他终于做出一件全然出于利他之心，而非为了满足自己的事，并因此迎来了至福的一刻。
浮士德随即死去，但由于这一刻源于善行，他没有立即坠入地狱，而是获得了接受审判的机会。
格蕾琴向圣母玛利亚恳求，请求由自己引领浮士德进入天国，圣母应允了她\cite{thomas2004introduction-faust}。


《魔圆》中，对《浮士德》的援引比比皆是。
除了前述的墙面涂鸦，魔女结界中的魔女文字也经常取自《浮士德》。
更重要的是，作品故事本身便充满了浮士德式的元素。
例如，“瓦尔普吉斯之夜”虽然在动画中被称作一个力量无匹的魔女，实际上却更像是一场由众多魔女共同参与的事件：
她们汇聚一处，纵情于毁灭的狂欢，恰如《浮士德》中描写的那般。
魔女以自身的绝望与疯狂篡写现实，由此筑成囚禁自己的结界。
这又令人联想起《浮士德》第一部结尾处陷入疯狂、身陷牢狱的格蕾琴。
对于浮士德而言，至福的一刻恰恰是通往地狱的入口，动画中的几个人物也经历了相似的逆转。
麻美刚刚沉浸在终于拥有朋友与同伴的幸福之中，片刻后便惨遭杀害；
杏子的父亲终于拥有了愿意聆听他布道的信众，正为此欣喜不已，却在得知这一切究竟是如何得来之后，胁迫家人与他一同自尽；
沙耶香也曾因自己拯救了恭介而感到幸福，却不曾料到，这份幸福正是她走向绝望、最终沦为魔女的起点。

更进一步，《魔圆》也可以视为对《浮士德》的一次改写。
丘比显然是梅菲斯特的化身：
它最初以可爱动物的形象示人，旋即又显露为一个可怖、强大而极具侵略性的存在，以实现愿望为饵，换取少女的灵魂。
正如梅菲斯特等待浮士德抵达幸福的顶点，再让他永远坠入地狱，丘比也利用魔法少女情感上的起伏，觊觎她们坠入绝望、化为魔女之际释放出的能量。

既然丘比最主要的目标是小圆，人们很容易顺势把她视为浮士德。
然而，这样的对应并不恰当。
本作更倾向于将小圆置于格蕾琴的位置；最直接的证据，便是她的魔女形态名为“救济的魔女\footnote{%
    译注：
    救济的魔女英文为 Kriemhild Gretchen，第二个单词与格蕾琴（Gretchen）的名字一致。
}”。
丘比在全剧大部分时间里反复诱使小圆订立契约，却屡屡失败，直至最后才终于得逞。
这与梅菲斯特最初几次企图腐化格蕾琴、令她倾心浮士德，却一再受挫的情形相互呼应。
最终，小圆许下愿望，要引导魔法少女逃离化为魔女的命运。
这恰与格蕾琴希望引领浮士德进入天堂的祈愿相映成趣。
小圆也由此成为救世主与守护者，其角色近似于《浮士德》中由“永恒的女性”所体现的圣母原则，而格蕾琴正与这一原则紧密相联\footnote{%
    译注：
    “永恒的女性”这一概念，指一种具有超越性、爱与宽恕力量的女性形象。
    而《浮士德》中的格蕾琴和海伦正是这一概念在尘世的两个象喻。
}\cite{obrien2012eternal-feminine}。


那么，倘若小圆不是浮士德，谁才是？
焰是个相当贴切的对象。
与浮士德一样，她同魔鬼订立契约，使时间倒流，试图补救那些她认为是自己造成的错误。
对焰而言，时间倒流确实是实际发生之事。
不过话说回来，浮士德也经历了类似穿越时空的冒险。
焰与小圆之间的深厚感情，以及她拯救小圆的愿望，也映照着浮士德对格蕾琴的感情。
她令时间停驻的能力，或许还暗合浮士德契约中那条关于“让此刻停留”的条件。
最终，她也像浮士德一样逐渐明白：
试图拨回时钟，并不能弥补过去，反而令一切愈发恶化。

然而，《魔圆》并非只是在复述《浮士德》，它也反过来颠覆了这部作品。
焰的愿望最终并不是一个错误，反而成为了打破循环的关键。
而小圆／格蕾琴为了获得引领他人进入天堂的力量，所求助的也不是圣母玛利亚或上帝，而是丘比／梅菲斯特。
这是因为，小圆既不真正对应《浮士德》中的任何人物，也根本不属于基督教传统。
她真正承担的，是另一套神话体系中的角色。

《魔圆》先披上一层典型魔法少女作品的“虚假”外衣，继而将其揭开，露出作品的“真实”面目。
与此相仿，它也先向观众呈现一种浮士德式、基督教式的表层诠释，借此遮蔽深处那个本质上属于佛教的故事。
最显而易见的对应在于，佛教的一项核心教义是有漏皆苦\footnote{%
    译注：
    “漏”指烦恼。所谓“苦”，不仅指肉体或心理痛苦，也包括无常所带来的“不圆满”。
    所以，“有漏皆苦”意指一切受烦恼、执著与业力支配的存在皆非圆满安稳。
}\cite{smart1998buddhism}，《魔圆》正是将这一命题转化成了自身的叙事机制：
所有愿望最终都会导向痛苦与绝望，情感每攀上一座高峰，便必然跌入与之相应的低谷。

本作还直接点明了另一个重要的佛教概念：业力。
业力的内涵极为复杂，不同佛教宗派对此也有截然不同的理解。
若仅作最宽泛的概括，佛教意义上的业，可以大致理解为因果相续：
人的行为如同播下种子，这些种子终将结成果报——或在今生，或在来世。
善业结善果，恶业结恶果，然而无论善恶，果报都会引发新的行为，新的行为又会种下新的因，由此使人始终困于业力的循环之中\cite{smart1998buddhism}。
魔法少女在幸福与绝望之间的循环也如出一辙，将她们一步步拖向魔女化的结局。

在佛教信仰中，业力的重负还会将众生束缚于生死轮回，使其一世又一世地重生，不断承担前世留下的业\cite{smart1998ingredients-indian-religion}。
这恰如第十集所描绘的那样：焰被困在周而复始的时间线中，始终无法脱身。
唯有开悟，也就是洞悉世界的真实本质，才能使人摆脱业力的循环。
而直至最后一次轮回，小圆才终于同时知晓两件事：
焰穿越时间的真相，也就是“轮回”，以及孵化者所作所为的完整真相，也就是“业力”的本质。
最后，瓦尔普吉斯之夜的形态酷似一朵莲花，而莲花正是开悟的象征\footnote{%
    这一象征在大乘佛教经典《法华经》中尤为突出。
    参见尼尼安·斯马特对日莲宗与日本民族主义的讨论\cite{smart1998nichiren-japanese-nationalism}。
}。
同时，作品中的齿轮意象也令人联想到不断转动着的法轮\cite{smart1998buddhism}。

如前所述，小圆真正近似的，其实是佛教中的观音菩萨——即日本的观音（Kannon）或中国的观音（Kuanyin）。
观音是一位几近进入涅槃的少女，却在即将跨过最后门槛时停下了脚步。
她超越时间与空间，向众生伸出援手，帮助他们获得开悟，最后才让自己进入涅槃。
这也解释了小圆为何会与圣母玛利亚产生联系：
观音与玛利亚在形象上的相似显而易见，也历来为论者所关注\footnote{%
    例如，尼尼安·斯马特曾指出观音与玛利亚形象之间的相似性。
}\cite{smart1993buddhism-christianity}。

由此看来，小圆先以菩萨之身跨越一切时空，救度所有众生。
继而越过最后的门槛，进入更高的境界。
她成为了自然之力本身，成为希望的化身。
她的个体意识随之消融，最终进入涅槃。

然而，正如后文将要看到的，她在抵达终点之前，稍稍绕了一段远路。
